\section{高阶紫外发散}
\label{sec:top}
从图~\ref{fig:oneloop-g}的一圈图出发，通过在其上添加胶子、正反夸克对或正反鬼场对，我们可以生成所有高阶圈图。由于准部分子算符的定义中 $z$ 分量被单独分离出来，圈动量 $\bar{l}$ 和 $l$ 的三维（3-D）与四维积分都可能导致紫外发散。在文献~\cite{Ishikawa:2017faj}中，我们分别针对高阶图的圈动量 3-D 与 4-D 积分引入了发散指数的改变量 $\Delta\omega_3$ 和 $\Delta\omega_4$，并证明：虽然并非必要，但只要所有相应高阶图都满足 $\Delta\omega_3\leq 0$ 与 $\Delta\omega_4\leq 0$，就足以保证准部分子算符可重正化。基于文献~\cite{Ishikawa:2017faj}导出的幂次计数规则，我们发现高阶下可能增大准胶子算符表观紫外发散的唯一情形是：添加一条两端都接在规范链上的胶子，此时 3-D 积分得到 $\Delta\omega_3=1$。把式~\eqref{eq:dec} 的分解应用于所加胶子的动量，即可直截了当地证明：任何圈数下经维数正规化的紫外发散都定域于时空之中，方式与准夸克算符相同~\cite{Ishikawa:2017faj}。于是我们发现 $\Delta\omega_3$ 实际上与研究紫外发散无关。又因为所有情形都有 $\Delta\omega_4\leq0$，所以只有图~\ref{fig:div-topo-g}中有限的几类高阶图拓扑具有紫外发散。

%%%%%%%%%%%%%%%%%%%%%%%%%%%%%%%
\begin{figure}[htb]
	\begin{center}
		\includegraphics[width=0.75\textwidth]{images/div-topo-g.pdf}
		\caption{\label{fig:div-topo-g}
			可能给准胶子算符带来紫外发散的四类图形拓扑。}
	\end{center}
\end{figure}
%%%%%%%%%%%%%%%%%%%%%%%%%%%%%%%

图~\ref{fig:div-topo-g}中拓扑 (a) 与 (c) 的泡泡表示单粒子不可约（1PI）图，二者都具有线性表观紫外发散。由于存在潜在的线性紫外发散，泡泡上再多接一条胶子的图可以产生对数紫外发散。产生对数发散的另一种可能是：一条胶子在泡泡之外但非常靠近泡泡处接到规范链上，如图~\ref{fig:div-topo-g} (b) 和 (d) 所示，其连接点用三角形标记。拓扑 (b) 与 (d) 的泡泡同时包含这里提到的两类对数发散图。

图~\ref{fig:div-topo-g}中的拓扑 (a) 与 (b) 与准夸克算符的相同，其发散可以用类似方式重正化。拓扑 (a) 图的线性发散可以通过一个整体因子消除，这相当于沿规范链运动的检验粒子的质量重正化~\cite{Polyakov:1980ca}；而由规范链端点引起的对数发散可以通过乘以 $Z_{wg}^{-1/2}$——检验粒子的“波函数”重正化——来消除~\cite{Dotsenko:1979wb}。拓扑 (b) 的图只有对数紫外发散，可由 QCD 重正化处理~\cite{Dotsenko:1979wb}。

图~\ref{fig:div-topo-g}中拓扑 (c) 与 (d) 图形的紫外发散不同于准夸克算符的情形，将在下面两节中分别研究。
